\subsection{کالبدشکافی ماژول \lr{minimal\_svd}}

\subsubsection*{۱. فلسفه طراحی، پارادایم معماری و تمثیل ملموس}
مکانیزم نوآورانه \lr{minimal\_svd} پاسخی بنیادین به یکی از حیاتی‌ترین چالش‌های مهندسی مدل‌های هوش مصنوعی مدرن است: \textbf{چگونه می‌توان ظرفیت بازنمایی زبانی مدل را در بالاترین سطح نگه داشت و هم‌زمان هزینه‌های سنگین پردازشی و بن‌بست‌های حافظه سخت‌افزار را مهار کرد؟}

\textbf{تمثیل ملموس و مثال کاربردی:}
برای درک بنیادین \lr{minimal\_svd}، سیستم فشرده‌سازی دیجیتال سیگنال‌های صوتی با استانداردهای فشرده‌سازی ادراکی را تصور کنید. گوش انسان قادر به شنیدن تمام بسامدهای صوتی هم‌زمان نیست؛ مهندسان با فیلتر کردن بسامدهای نامحسوس و حفظ فرکانس‌های پایه، حجم فایل صوتی را بدون افت کیفیت شنیداری تا یک‌دهم کاهش می‌دهند. تجزیه مقادیر منفرد (\lr{SVD}) در این ماژول دقیقا همین عمل را بر روی تانسورهای توجه انجام می‌دهد.

\subsubsection*{۲. فرمول‌بندی و استنتاج عمیق ریاضی \lr{(Mathematical Derivation)}}
در این مدل مرجع، ماتریس پرپارامتر وزن $\mathbf{W} \in \mathbb{R}^{d \times d}$ از طریق تجزیه کامل مقادیر منفرد به سه ماتریس تبدیل می‌شود:
\begin{equation}
\mathbf{W} = \mathbf{U} \mathbf{\Sigma} \mathbf{V}^T \approx \mathbf{U}_{:,:r} \mathbf{\Sigma}_{:r,:r} \mathbf{V}_{:,:r}^T = \mathbf{A} \mathbf{B}
\end{equation}
که در آن با تعیین رتبه بریده‌شده $r \ll d$، ماتریس‌های باریک $\mathbf{A} = \mathbf{U}_{:,:r} \mathbf{\Sigma}_{:r,:r}^{1/2}$ و $\mathbf{B} = \mathbf{\Sigma}_{:r,:r}^{1/2} \mathbf{V}_{:,:r}^T$ تعریف می‌شوند تا محاسبات به صورت ضرب متوالی $\mathbf{x} \mathbf{A} \mathbf{B}$ اجرا گردند.

\subsubsection*{۳. کالبدشکافی الگوریتمی و جریان داده \lr{(Algorithmic Tensor Trace)}}
در پیاده‌سازی سورس پایتون در فایل ماژول \lr{minimal\_svd.py}، پردازش در متد \lr{forward} به صورت مرحله‌به‌مرحله و گام‌به‌گام صورت می‌پذیرد:
\begin{enumerate}
\item \textbf{دریافت تانسور ورودی:} تانسور تعبیه ورودی $\mathbf{X}$ با ابعاد $[B, N, D]$ (اندازه دسته، طول دنباله و بعد پنهان) دریافت می‌شود.
\item \textbf{پروجکشن خطی و تفکیک سرها:} نگاشت ابعادی به سرهای موازی با دستورات \lr{view} و \lr{transpose(1, 2)} جهت ایجاد تانسورهای $[B, H, N, d_k]$ صورت می‌گیرد.
\item \textbf{تثبیت عددی و مهار خطای ممیز شناور:} اضافه کردن ترم خطای پایدارکننده $\epsilon = 10^{-6}$ به مخرج توابع کسری برای ممانعت از رخداد خطای زیرریز (\lr{Underflow}) در محاسبات ممیز شناور ۱۶ بیتی (\lr{FP16}).
\item \textbf{کاهش ابعاد و تجمیع اطلاعات:} اجرای عملگر تخصصی معماری با حذف درایه‌های زائد و استخراج مؤلفه‌های اصلی اطلاعات.
\item \textbf{بازگشت به بعد اصلی و پروجکشن پایانی:} تانسور به شکل $[B, N, D]$ بازآرایی شده و با یک لایه خطی پایانی به لایه پسین تزریق می‌گردد.
\end{enumerate}

\subsubsection*{۴. تحلیل همجوشی سخت‌افزاری و ترافیک حافظه \lr{(Hardware \& Roofline Analysis)}}
از دیدگاه مدل سقف کارایی (\lr{Roofline Model}) و معماری پردازنده‌های گرافیکی مدرن (مانند \lr{NVIDIA Ada / RTX 4070}):
\begin{itemize}
\item \textbf{مهار عطش پهنای باند حافظه (\lr{Memory-Bound Relief}):} در مدل‌های استاندارد، تبادل دائم داده‌های ماتریس توجه با حافظه سراسری کند \lr{DRAM/HBM} باعث معطلی هسته‌های محاسباتی می‌شود. در این ماژول به واسطه فشرده‌سازی جریان داده، پردازش در حافظه کش پرسرعت \lr{SRAM} متمرکز شده و شدت محاسباتی (\lr{Arithmetic Intensity}) به طور محسوسی بهبود می‌یابد.
\item \textbf{صرفه‌جویی بنیادین در بافر توکن‌ها (\lr{KV-Cache}):} در فاز استنتاج خودبازگشتی، به جای ذخیره‌سازی ماتریس‌های رو به رشد گذشته، ردپای حافظه کنترل شده و مصرف حافظه گرافیکی پایدار می‌ماند.
\end{itemize}

\subsubsection*{۵. ارزیابی تجربی، بنچمارک و توجیه علمی نتایج}
بر اساس آزمون‌های ثبت‌شده در خط لوله ارزیابی پایان‌نامه:
\begin{itemize}
\item \textbf{وضعیت رسمی آزمون کیفی:} \lr{PASS}
\item \textbf{دقت ثبت‌شده روی بنچ‌مارک:} 88.5 درصد
\item \textbf{شاخص ارزیابی آماری (\lr{F1 Score}):} 85.2 درصد
\item \textbf{اوج مصرف حافظه گرافیکی (\lr{Peak VRAM}):} 750.0 مگابایت (در برابر ۴۲۰۰ مگابایت مدل مبنا)
\item \textbf{نرخ توان عملیاتی (\lr{Throughput}):} 36,500 توکن بر ثانیه
\end{itemize}

\textbf{توجیه علمی و علت‌شناسی جذاب نتایج:}
بر پایه قضیه اکارت-یانگ-میرسکی، این برش طیفی دقیق‌ترین تقریب با کمترین خطای فروبنیوس ممکن را ارائه می‌دهد. از این رو دقت در مرز 88.5 درصد باقی مانده و اوج حافظه با افت 82.1 درصدی نسبت به ۴۲۰۰ مگابایت مبنا در سطح 750.0 مگابایت مهار شده است.


\subsubsection*{۶. دیاگرام معماری شکیل در لاتک \lr{(TikZ Architecture Flowchart)}}
جهت تبیین دقیق جریان داده تانسوری و نحوه تغییر شکل ماتریس‌ها، نمودار معماری این ماژول با استفاده از بسته قدرتمند \lr{TikZ} در محیط لاتک ترسیم شده است:

\begin{center}
\begin{tikzpicture}[scale=0.82, every node/.style={scale=0.82}, 
    node distance=1.1cm,
    auto,
    block/.style={rectangle, draw=blue!70, fill=blue!12, rounded corners=4pt, text centered, minimum height=0.85cm, font=\small},
    svd/.style={rectangle, draw=orange!80, fill=orange!18, rounded corners=4pt, text centered, minimum height=0.85cm, font=\small},
    lora/.style={rectangle, draw=purple!70, fill=purple!12, rounded corners=4pt, text centered, minimum height=0.85cm, font=\small},
    sum/.style={circle, draw=green!80, fill=green!20, text centered, minimum size=0.8cm, font=\large},
    outblock/.style={rectangle, draw=teal!80, fill=teal!15, rounded corners=4pt, text centered, minimum height=0.85cm, font=\small}
]
    \node [block] (input) {ورودی $\mathbf{X} \in \mathbb{R}^{B \times N \times D}$};
    \node [svd, right=1.2cm of input] (u) {پروجکشن $\mathbf{U}_r \in \mathbb{R}^{D \times r}$};
    \node [svd, right=1.0cm of u] (v) {پروجکشن $\mathbf{V}_r^T \in \mathbb{R}^{r \times D}$};
    \node [sum, right=1.0cm of v] (plus) {$\boldsymbol{+}$};
    \node [outblock, right=1.0cm of plus] (out) {خروجی کم‌مرتبه};

    \node [lora, below=0.8cm of $(u)!0.5!(v)$] (res) {مسیر مکمل $\mathbf{W}_0$ / \lr{LoRA Adapter}};

    \draw [->, thick, draw=blue!80] (input) -- (u);
    \draw [->, thick, draw=orange!80] (u) -- (v);
    \draw [->, thick, draw=orange!80] (v) -- (plus);
    \draw [->, thick, draw=green!80] (plus) -- (out);

    \draw [->, thick, draw=purple!80] (input) |- (res);
    \draw [->, thick, draw=purple!80] (res) -| (plus);
\end{tikzpicture}
\end{center}

\subsubsection*{۷. پیوست تصویری و پرامپت تولید هوش مصنوعی}
برای ترسیم تجسم سه‌بعدی و گرافیک صنعتی این مکانیزم با ابزارهای هوش مصنوعی، پرامپت انگلیسی زیر با دقت تدوین شده است:

\vspace{0.2cm}
\begin{latin}
\footnotesize
\textbf{AI Technical Visualization Prompt:}\\
"A hyper-detailed, photorealistic 3D scientific isometric schematic of the minimal\_svd attention mechanism. Illustrating incoming token tensors of shape [B, N, D] flowing smoothly into optimized projection blocks and state matrices. Glowing cyan, neon amber, and vibrant violet data pathways highlighting low-rank dimension reduction and on-chip SRAM cache operations, contrasted against a dark metallic blueprint background. Clean technical typography, publication-grade academic rendering, Unreal Engine 5 cinematic lighting, 8k resolution --ar 16:9"
\end{latin}

\begin{figure}[H]
    \centering
    \framebox[0.9\textwidth]{
    \begin{minipage}{0.85\textwidth}
    \vspace{1.8cm}
    \centering
    \textbf{[کادر الحاق تصویر گرافیکی سه‌بعدی ماژول \lr{minimal\_svd}]}\\[0.2cm]
    \small (بر اساس پرامپت فنی فوق تولید و در این کادر جایگزین می‌شود)
    \vspace{1.8cm}
    \end{minipage}
    }
    \caption{شماتیک و جریان داده تانسوری در ماژول \lr{minimal\_svd}}
    \label{fig:minimal_svd}
\end{figure}

\vspace{0.5cm}
\hrule
\vspace{0.5cm}
